% \chapter{Imagining Biophilic Interaction in Indoor Environments}

% Having examined how an affective lens can help HBI understand existing indoor experience, the dissertation shifts towards design. This transition raises a further question: how can experiential knowledge inform future indoor environments without being converted into requirements for producing predetermined emotional states?

% To address this question, the thesis engages architecture, spatial-design, and related experts. Their perspectives are important because the design of indoor environments involves knowledge of spatial organisation, sensory conditions, materiality, atmosphere, building use, and architectural practice. The expert studies bring these forms of knowledge into dialogue with the experiential concerns of HCI and HBI.

% The first study consists of semi-structured interviews examining expert imaginaries of multisensory biophilic design.\sidenote{I use \emph{imaginaries} to refer to future-oriented accounts of what indoor environments could become. The term concerns design possibility rather than prediction.} Rather than asking experts to evaluate an established system, the interviews invite them to articulate possibilities, tensions, and directions for future indoor environments.

% Biophilic design provides a productive context for this inquiry because it concerns how relations with nature can be incorporated into built environments. It is frequently represented through plants, daylight, natural materials, water, views, or organic forms. These strategies are important, but they do not encompass the full experiential potential of relations with nature.

% Nature-related experience can also involve sound, movement, temporal variation, atmosphere, ambiguity, sensory correspondence, and changing environmental conditions. These qualities connect to the affective account developed through the occupant studies. They provide ways of reasoning about experience without assuming that design should generate a single desirable emotional state.

% The interviews identify dimensions of multisensory biophilic design that remain under-explored within technologically mediated indoor environments. They direct attention beyond the visual presence of nature towards environmental processes, sensory relationships, social experience, spatial context, and the ways interactive technologies may support or complicate these relations.

% The interviews establish the relevant dimensions, but they do not fully demonstrate how experts would express them through design. The dissertation therefore proceeds to expert design sessions in which participants articulate, visualise, and discuss possible multisensory biophilic environments. These sessions examine how future-oriented ideas are translated into spatial and interactive propositions.

% Together, the interviews and design sessions produce a biophilic design framework for HBI, with broader implications for HCI. The framework organises the expert findings into design approaches for technologically mediated indoor environments. It does not prescribe a fixed set of natural elements or system features. Its purpose is to expand the design vocabulary of HBI by connecting affective, multisensory, spatial, and nature-related experience.

\chapter{Expanding the Experiential Scope of Biophilic Interaction}
\label{ch:biophilic-futures}

\begin{synopsis}
\synopsislead{What can biophilic architecture contribute to the design}{of interactive experiences in future indoor environments? The preceding chapter established Affective Interaction as a methodological lens for investigating the complexity of lived indoor experience. This chapter turns from investigating existing experience towards considering how future interactive environments might support richer forms of human--nature connection.}

Biophilic design provides one direction for this inquiry. Architecture offers a long history of shaping connections with nature through sensory-rich spaces, natural materials, spatial organisation, environmental processes, and ecological relations. Within technologically mediated indoor environments, however, biophilic interaction remains comparatively limited, with much of its potential for multisensory and multidimensional experience still underexplored.

This chapter examines what interaction design can learn from this architectural knowledge through semi-structured interviews with 13 experts. The analysis develops eight themes reflecting expert imaginaries of future biophilic environments and five design opportunities: designing for Emergent Natural Processes, Occupant Body Sensitivity, Biophilic (Dis)Comfort, Biophilic Equity, and Relational Encounters with Nature. Together, these findings position biophilic interaction as multisensory and interpretive, while extending its scope towards more-than-human and justice-oriented considerations.

Within the trajectory of the dissertation, this chapter develops \rqref{biophilic} by broadening the experiential possibilities available to the design of future interactive indoor environments. It establishes how architectural knowledge can extend biophilic interaction beyond predominantly visual or restorative representations of nature towards richer sensory, ecological, and relational forms of experience. These possibilities provide the basis for the following chapter, which examines how they can be translated into a more structured interaction design framework.
\end{synopsis}

\chaptersource{%
This chapter is based on: Shruti Rao, Judith Good, and Hamed Alavi.
\emph{Designing Multisensory Biophilic Futures: Exploring the Potential of Interaction Design to Deepen Human Connections with Nature in Indoor Environments}.
Presented at DIS 2025, Funchal, Portugal.%
}


\import{papers/03-multisensory-biophilic/}{thesis-body.tex}
\FloatBarrier
\clearpage

\begin{chapterclosing}{What this chapter establishes}

This chapter establishes multisensory biophilic design as an approach for broadening the experiential scope of interactive indoor environments. The expert accounts show that meaningful connections with nature can extend beyond visual representations or isolated natural features to include sensory richness, environmental processes, bodily experience, ecological relations, and forms of interpretation that develop through encounters with nature.

The eight themes articulate this broader understanding of biophilic experience, while the five design opportunities translate it into directions for future interactive environments: Emergent Natural Processes, Occupant Body Sensitivity, Biophilic (Dis)Comfort, Biophilic Equity, and Relational Encounters with Nature. Across these directions, multisensory and multidimensional experience provides an underlying orientation, while more-than-human and justice-oriented concerns extend biophilic interaction beyond individual wellbeing alone. :contentReference[oaicite:3]{index=3}

Within the dissertation, this chapter establishes the experiential and conceptual scope of the biophilic design direction. Architectural expertise expands the design vocabulary available to HBI by foregrounding forms of human--nature connection that are sensory, interpretive, ecological, and relational. The resulting themes and opportunities identify a broad design territory for future interactive indoor environments, while leaving open how its constituent experiential qualities can be organised and translated into more systematic interaction design concepts.

\chaptercarryforward{The following chapter develops this design direction through expert design sessions, moving from a broad account of multisensory biophilic possibilities towards a structured framework that specifies what can be configured, why particular experiential qualities matter, and how biophilic interaction might occur within indoor environments.}
\end{chapterclosing}